%قبل از شروع مسئله ترجیحاً چیزی ننویسید. اما اگر خواستید بنویسید، در قالب لیستی (فقرات و شمارش) بنویسید.
%برای اضافه کردن نمره، بعدی از شروع‌مسئله، [x] را اضافه کنید و به جای x  نمره مورد نظر را قرار دهید. برای زیرمسئله‌ها هم به همین صورت است.
%\شروع‌مسئله[10]
\شروع‌مسئله
تورنمنت یک گراف کامل جهت‌دار است. به دو سوال زیر پاسخ دهید.
% صورت اصلی مسئله را در این‌جا بنویسید.
\begin{enumerate}
\item
ثابت کنید در هر تورنمت راسی وجود دارد که عمق درخت $BFS$ مربوط به آن راس حداکثر 2 است.
\item
باتوجه به قسمت قبل ثابت کنید اگر هیچکدام از رئوس درجه‌ی خروجیشان $n-1$ نباشد، حداقل دو راس با چنین خاصیتی داریم.
\end{enumerate}

\solution{
\begin{enumerate}
\item
دقت کنید حکم داده شده در قسمت اول معادل با این است که ثابت کنیم در یک تورنمت راسی وجود دارد که با حداکثر دو یال به تمامی رئوس مسیر دارد.(چنین راسی را شاه می نامیم)برای اثبات از استقرا استفاده می کنیم.پایه استقرا $n=1$ و $n=2$ است که حکم برای این دو پایه بدیهی است.برای گام استقرا ابتدا راس دلخواه $a$ را حذف کنید.روشن است که یک تورنمت با $n-1$ راس باقی می ماند.طبق فرض استقرا این تورنمت شاهی دارد که آن را $b$ می نامیم.حال راس $a$ را به گراف بازگردانید.اگر $b$ به $a$ یال داشته باشد که کار تمام است و $b$ همچنان شاه است.نتیجتا فرض می کنیم $a$ به $b$ یال دارد.حال رئوسی را که $b$ با آن ها یال دارد را درنظر بگیرید.اگر هریک از این رئوس به $a$ یال داشته باشند آنگاه $b$ با دو یال به $a$ مسیر دارد که این یعنی مجددا $b$ شاه می ماند.فرض کنید $a$ به تمامی این رئوس یال داشته باشد.با کمی دقت می توان فهمید که در این صورت $a$ از نظر سطح دسترسی مانند $b$ شده است(تمامی رئوسی که $b$ به آنها یال داشت $a$ هم یال دارد پس تمامی رئوسی که $b$ با دئ یال به آنها مسیر داشت $a$ هم با دو یال به آنها مسیر دارد) با این تفاوت که $a$ به خود $b$ هم یال دارد نتیجتا $a$ جدید تورنمت است.
\item
چون راسی نیست که درجه خروجی اش $n-1$ باشد نتیجتا سطح دوم درخت $BFS$ خالی نیست.از طرفی می دانیم که تورنمت قطعا یک شاه را دارد.این شاه و تمام رئوسی که به آن ها یال دارد را حذف کنید.گراف باقی مانده مجددا یک تورنمت است و شاه مخصوص به خودش را دارد.همچنین بدیهی است که شاه این تورنمت کوچکتر باید به شاه تورنمت بزرگتر یال داشته باشد(وگرنه حذف میشد) پس شاه تورنمت کوچکتر به تمام رئوسی که شاه اول با یک یال به آنها مسیر دارد با دو یال مسیر دارد.از تمامی این جملات جمع بندی میکنیم که این شاه دوم هم یک شاه برای کل تورنمت است.
\end{enumerate}
}



%%%%%%%%%%%%%%%%%%%%%%%%%%%%%%%%%%%%%%%%%%%%%%%%%%%%%
% برخی دستورات مفید برای تمرین‌ها. %
%%%%%%%%%%%%%%%%%%%%%%%%%%%%%%%%%%%%%%%%%%%%%%%%%%%%%

% فرمول‌های نمونه:
%$f(x) = \Theta(x)$
%$T(x, y) \leq \log(x) + 2^y + \binom{2n}{n}$
% {\tt code\_function}


% لیست موارد بدون شماره‌گذاری:
%\شروع{فقرات}
%    \فقره مورد اول 
%        \شروع{فقرات}
%            \فقره مورد اول داخلی 
%                \شروع{فقرات}
%                    \فقره اولین مورد خیلی داخلی 
%                    \فقره دومین مورد خیلی داخلی 
%                \پایان{فقرات}
%            \فقره مورد دوم داخلی 
%        \پایان{فقرات}
%    \فقره مورد دوم 
%\پایان{فقرات}

% لیست موارد شماره‌گذاری شده:
%\شروع{شمارش}
%    \فقره یک مورد 
%    \فقره یک مورد دیگر 
%    \فقره باز هم یک مورد دیگر
%\پایان{شمارش}

% برای نوشتن یک فرمول چند خطی که یک شماره بخوره. برای حالت‌های دیگر به سایت زیر مراجعه کنید:
% https://www.overleaf.com/learn/latex/Aligning_equations_with_amsmath
%\begin{align} 
%    \begin{split}
%        (x+y)^3 &= (x+y)^2(x+y) \\
%                &= (x^2+2xy+y^2)(x+y) \\
%                &= (x^3+2x^2y+xy^2) + (x^2y+2xy^2+y^3) \\
%                &= x^3+3x^2y+3xy^2+y^3
%    \end{split}                                 
%\end{align}

% نوشتن ماتریس و آرایه دوبعدی
%\begin{align}
%    A = 
%    \begin{bmatrix}
%        A_{11} & A_{21} \\
%        A_{21} & A_{22}
%    \end{bmatrix}
%\end{align}

% اضافه کردن تصویر با شرح. برای اضافه کردن تصویر بدون شرح، فقط دستور درج‌تصویر را استفاده کنید.
%\شروع{شکل}[h!]
%\تنظیم‌ازوسط
%  \درج‌تصویر[پهنا=0.5\پهنای‌متن]{img.jpg}
%\شرح{شرح تصویر}
%\برچسب{تصویر:۱}
%\پایان{شکل}
%\پایان{زیرمسائل}

%کد زیر را ببینید.
%%\lstinputlisting{source_filename.py}
%\begin{latin}
%\begin{lstlisting}
%# Import the modules
%import sys
%import random
%
%ans = True
%
%while ans:
%question = raw_input("Ask the magic 8 ball a question: (press enter to quit) ")
%
%answers = random.randint(1,8)
%
%if question == "":
%sys.exit()
%
%elif answers == 1:
%print "It is certain"
%
%elif answers == 2:
%print "Outlook good"
%
%elif answers == 3:
%print "You may rely on it"
%
%elif answers == 4:
%print "Ask again later"
%
%elif answers == 5:
%print "Concentrate and ask again"
%
%elif answers == 6:
%print "Reply hazy, try again"
%
%elif answers == 7:
%print "My reply is no"
%
%elif answers == 8:
%print "My sources say no"
%\end{lstlisting}
%\end{latin}
